\chapter{基于矩阵QR分解技术的编队量测配置}

\section{引言}
随着无人机集群技术向大规模、低成本及实战化方向演进，复杂动态环境下的状态感知与获取已成为制约系统性能的关键瓶颈。尽管分布式协同控制理论已取得显著进展，但现有的主流算法通常建立在个体具备全维、高精度状态信息的理想化假设之上。然而，在实际工程应用中，受限于物理载荷约束与强对抗环境，这一假设往往难以维系。这种理论假设与工程现实之间的鸿沟，引出了当前无人机集群状态感知领域面临的三大核心挑战：

\textbf{1）全维高精度传感的物理约束与成本博弈：} 
传统全维状态感知方案严重依赖于RTK-GNSS、激光雷达或高精度惯导系统。然而，在大规模集群“低成本、可消耗”的战术背景下，单机载荷能力与制造成本的严格限制，使得为每个节点配置全套高精度传感器成为不可行的工程选项。这导致集群中绝大多数节点仅能依靠低精度传感器获取非全维状态信息，无法独立构建全局坐标系下的完整状态向量，从而限制了集群的整体感知精度。

\textbf{2）拒止环境下的定位失效与单一量测局限：} 
在复杂场景中，外部卫星导航信号面临衰减、欺骗乃至完全中断的风险。一旦失去外部绝对定位信息，基于全维状态反馈的传统编队控制闭环将因输入缺失而失效，进而引发队形发散或碰撞。虽然相对距离量测具有不依赖外部设施且抗干扰能力强的优势，但仅凭距离信息存在平移与旋转的模糊性。此外，尽管采用纯方位角量测可规避对绝对位置的依赖，但其引入了新的几何约束难题。文献\parencite{CHEN2022110605}指出，二维空间中仅依赖方位角的系统，其可观测性高度耦合于智能体间的相对几何构型；一旦系统演化至共线等退化构型，整个编队将陷入不可定位的奇异状态。这表明单一模态的相对量测在复杂动态环境中难以独立维持系统的可定位性。

\textbf{3）异构量测数据的融合与可观测性量化难题：} 为应对上述挑战，“异构传感配置”策略成为一种有效的解决方案。然而，这种混合模式引发了关于系统可观测性判定与保障的深层理论问题。现有研究如文献\parencite{1025347369.nh}指出，编队稳定性与通信拓扑的刚性密切相关。此外，文献\parencite{10197184}进一步从图论视角探讨了多智能体系统的刚性与双刚性条件，为基于距离量测的网络结构提供了定性分析基础。然而，上述理论框架多侧重于定性判断，对于如何在给定通信拓扑下，量化不同锚点配置对系统整体可观测性的贡献权重，并据此优化传感器布局，尚缺乏一套普适且高效的代数分析方法。若缺乏定量的优化指导，即便满足拓扑刚性条件，编队仍可能存在不可观测的“柔性模态”。

针对上述量测模式的局限性与现有理论工具在量化分析上的不足，本章旨在寻求一种基于代数图论与矩阵分析的新型解决方案。我们将突破单纯依赖几何刚性判定的传统范式，引入矩阵QR分解技术对系统的观测矩阵进行数值分析。本章致力于建立一套可量化的传感器配置理论框架，通过精确评估不同节点作为锚点的“信息价值”，最终确定满足系统“弱可观测性”要求的最优或次优锚点部署策略。这项工作不仅是对现有刚性理论的有力补充，更是解决异构集群协同感知这一核心难题的关键一步，为后续设计鲁棒的分布式状态观测器奠定坚实的理论基础。

\section{非全维量测环境下的编队感知建模}

\subsection{非结构化场景下典型传感器的局限性分析}

在无人机集群的实际工程应用中，作业环境往往具有高度的非结构化特征。与实验室内具备高精度运动捕捉系统的理想环境不同，此类场景对机载传感器的感知能力构成了严峻挑战。现有的典型定位传感器在非结构化场景中面临着不同程度的失效风险，具体局限性分析如下：

\textbf{1) GNSS的易受扰性与盲区问题\cite{2022A, XDDH202506001}：}
全球导航卫星系统是目前获取室外绝对位置信息的主要手段，但其信号功率极其微弱，极易受外部环境干扰。在城市峡谷、茂密林区或山谷等遮挡环境下，多径效应会导致定位精度急剧下降并产生显著漂移。
而在军事对抗等强电磁干扰场景中，GNSS信号不仅可能被噪声压制，更可能遭受欺骗式干扰，导致无人机解算出错误的全局坐标而彻底迷失方向。
现有研究指出，即便采用惯性导航与卫星信号的深组合技术，一旦卫星信号完全中断，系统将退化为纯惯性导航模式，无法从根本上解决持续定位问题。这意味着在关键任务阶段，集群必须具备在无卫星信号参考条件下的自主生存能力。

\textbf{2) 视觉与激光传感器的环境依赖性与计算负荷\cite{Lu2018ASO,8517639}：}
视觉里程计和激光雷达虽能在无卫星环境下提供相对位姿估计，但其性能高度依赖环境特征的丰富度。视觉传感器在光照剧烈变化、弱纹理区域或烟雾粉尘遮蔽的非结构化环境中，特征提取与匹配极易失效从而导致跟踪丢失。
激光雷达虽然具有主动测距优势，但受限于尺寸、重量、功耗及成本，高性能三维激光雷达难以在大规模微型无人机集群中普及，且在特定角度下存在扫描盲区。
此外，这两类传感器产生的高维数据流对机载嵌入式芯片的实时处理能力提出了极高要求，难以满足大规模集群对低延迟与低功耗的工程约束。

\textbf{3) 惯性测量单元（IMU）的积分漂移累积\cite{Sahawneh2008DevelopmentAC}：}
微机电系统惯性测量单元（MEMS IMU）因其高频输出和抗干扰能力成为导航核心，但低成本商用器件存在显著的零偏不稳定性、尺度因子非线性及温度漂移误差。
相关研究表明，未经高阶校准的低成本惯性器件，其角随机游走和速度随机游走误差会随时间快速累积。在缺乏外部绝对位置校正的情况下，仅依靠惯性测量单元进行二次积分推算位置，误差将呈时间的多项式规律发散。这使得低成本惯性导航仅能维持秒级的高精度航位推算，无法独立支撑长时间的编队飞行任务，必须依赖外部量测进行在线校准。

\textbf{4) 气压计与测距传感器的局部分维特性\cite{ZJXH201510001033,7346960}：}
气压计作为一种低成本传感器能提供相对稳定的高度信息，但其测量易受近地气流扰动和温度变化影响，且仅能提供垂直方向的一维状态，无法解算水平位置。
超宽带技术（UWB）虽能在室内等非视距环境下提供高精度的节点间相对距离量测，但其定位精度受锚点几何分布及多径效应影响显著，且通常仅能提供相对拓扑约束而非绝对地理坐标。这两类传感器虽具备较强的鲁棒性，但均属于非全维量测，单独使用无法构建完整的六自由度状态估计。

综上所述，在非结构化复杂场景中，期望每一架无人机时刻拥有全维且高精度的绝对定位信息并不符合工程实际。单一传感器的物理局限性决定了其在极端环境下的不可靠性，而传统依赖全量测信息的控制理论在此类场景下往往失效。为了提高无人机集群在真实环境中的工程可行性与生存能力，必须突破传统思维，研究如何在传感器配置受限、仅能获取非全维或部分相对量测信息的条件下，通过多源异构信息融合与集群协同观测，弥补单一传感器的缺陷以实现系统的可观测性增强。这正是本章建立非全维量测模型、推导可观测性条件并研究传感器优化配置的核心动机与现实意义所在，旨在为资源受限下的无人机集群稳健导航提供理论支撑。

\subsection{异构量测系统的数学描述}
设$ i \in \{1,2,\dots,N\} $为集群中各无人机的索引，其中$ N>3 $表示集群规模。为兼顾无人机全局位置的可观测性分析与期望速度指令的生成，本节基于单积分器运动学模型(\ref{位置导数等于速度})，将第 $i$ 架无人机的动态演化过程描述为：
\begin{equation}
	\dot{\mathbf{p}}_i = \mathbf{v}_i^*,
	\label{一阶动力学}
\end{equation}
式中，$ {\mathbf{p}}_i = [x_i;y_i;z_i]$ 与 $\mathbf{v}_i^* = [{v}_{x,i}^* ; {v}_{y,i}^*; {v}_{z,i}^*] $ 分别表示无人机 $i$ 在以子群中心为原点的全局坐标系下的位置向量与期望速度指令 。

针对非全维位置传感特性，无人机 $i$ 的绝对位置量测方程可建模为：
\begin{equation}
	\tilde{\mathbf{z}}_i = \tilde{C}_i \mathbf{p}_i,
	\label{位置量测输出}
\end{equation}
此处，$\tilde{\mathbf{z}}_i$ 表示位置传感器的实际量测输出 。$\tilde{C}_i \in \mathbb{R}^{3 \times 3}$ 为对角量测矩阵，其主对角线元素取值为 $0$ 或 $1$，物理意义为对应维度的绝对位置是否可测 。例如，若无人机 $i$ 仅搭载气压计测量高度或水平定位模块失效，其量测矩阵则退化为 $\tilde{C}_i=\mathrm{diag}(0,0,1)$ 。

对于无人机间的测距传感器，其量测为$ ||{\mathbf{p}}_i-{\mathbf{p}}_j||,j \in \mathcal{N}_i $，为便于后续描述与分析，定义变量$ \beta_{ij} = \frac{1}{2} ||{\mathbf{p}}_i-{\mathbf{p}}_j||^2 $ ，并按照边索引的方式定义非线性量测方程：
\begin{equation}
	\bar{\mathbf{z}}_i=[\cdots; \bar{z}_{im}; \cdots], \quad \bar{z}_{im}=|h_{im}|\beta_{ij}, \quad
	m \in \{1,2,\dots,M\},
	\label{距离量测方程}
\end{equation}
式中$ |h_{im}| \in {0,1} $为关联矩阵 $ \mathcal{H} $ 的元素，表征网络拓扑中的第 $m$ 条通信链路是否与节点 $i$ 直接相连 。若 $|h_{im}|=0$，则意味着该边对应的距离信息对无人机 $i$ 是不可见的 。综合上述绝对与相对量测信息，无人机 $i$ 的全源量测向量可堆叠集成如下：
\begin{equation}
	\mathbf{z}_i = [\bar{\mathbf{z}}_i;\tilde{\mathbf{z}}_i] = 
	\left[
		\underbrace{
			\left[
			\cdots; \frac{1}{2} ||{\mathbf{p}}_i-{\mathbf{p}}_j||^2;\cdots
			\right]
		}_{\text{距离量测输出}};
		\underbrace{
			\left[
			\tilde{C}_i \mathbf{p}_i
			\right]
		}_{\text{位置量测输出}}
	\right],
\end{equation}
对于非全连通的测距传感器网络，无人机$i$无法测量集群内某些无人机对应的距离，因此在量测方程(\ref{距离量测方程})中通过系数$ |h_{im}| = 0 $来区分。

根据无人机的位置量测条件，可将编队中的无人机分为三类：

\textbf{1）锚点无人机：}具有全维的位置量测，对应的位置输出矩阵为$\tilde{C}_i = I_3$；

\textbf{2）受限锚点无人机：}具有非全维的位置量测，对应的位置输出矩阵$\tilde{C}_i \neq I_3$且$\tilde{C}_i \neq \mathbf{0}_{3 \times 3}$；

\textbf{3）非锚点无人机：}仅具有相对位置量测，不具备位置量测，对应的位置输出矩阵$\tilde{C}_i=\mathbf{0}_{3 \times 3}$。

整个集群的异构量测架构本质上由各节点量测矩阵 $\tilde{C}_i$ 的主对角元素共同定义 。由于协同编队控制律的生成高度依赖于准确的全局位置与相对状态反馈，对于非锚点或受限锚点无人机，必须设计基于异构量测信息的分布式全局位置观测器 。为从宏观层面分析整个编队系统的观测特性，定义集群全局位置状态向量为：
\begin{equation}
	{\mathbf{p}} = [\cdots;{\mathbf{p}}_i;\cdots].
\end{equation}
相应地，全局量测输出向量 $\mathbf{z}$ 堆叠了所有机间距离与绝对位置信息：
\begin{equation}
	\mathbf{z} = [\bar{\mathbf{z}};\tilde{\mathbf{z}}] = 
	\left[
		\underbrace{
			\left[
			\cdots; \frac{1}{2} ||{\mathbf{p}}_i-{\mathbf{p}}_j||^2;\cdots
			\right]
		}_{\text{距离量测输出}};
		\underbrace{
			\left[
			\cdots ; 
			\tilde{C}_i \mathbf{p}_i ;
			\cdots 
			\right]
		}_{\text{位置量测输出}}
	\right],
\end{equation}
式中$ \bar{\mathbf{z}} = [\cdots; \beta_{ij}; \cdots] $ 表示编队无人机间的距离量测，$ \tilde{\mathbf{z}} = \left[\cdots; \tilde{C}_i \mathbf{p}_i ; \cdots \right]$ 表示编队的位置量测。然后可以定义整个编队系统的输出方程为：
\begin{equation}
	\mathbf{z} = h(\mathbf{p}).
\end{equation}
参考李导数的定义，无人机编队系统在位置$ \mathbf{p} $处的可观测性矩阵可推导为：
\begin{equation}
	\mathcal{O}_{\mathbf{p}} = 
	\begin{bmatrix}
		\mathcal{R}_{\mathbf{p}}  \\
		\tilde{C} 
	\end{bmatrix},
	\label{可观测性矩阵}
\end{equation}
式中$ \mathcal{R}_{\mathbf{p}}$表示编队在$ \mathbf{p} $处的刚性矩阵，$\tilde{C} = \text{diag}(\cdots, \tilde{C}_i, \cdots)$是各无人机位置输出矩阵构成的分块对角矩阵，可以写成：
\begin{equation}
	\tilde{C} = 
	\begin{bmatrix}
		\tilde{C}& & & & \\
		& \ddots & & & \\
		& & \tilde{C}_i & & \\
		& & & \ddots & \\
		& & & & \tilde{C}_N
	\end{bmatrix}
	=
	\begin{bmatrix}
		\tilde{c}& & & & \\
		& \ddots & & & \\
		& & \tilde{c}_s & & \\
		& & & \ddots & \\
		& & & & \tilde{c}_{Nn}
	\end{bmatrix}.
\end{equation}
一旦对角线元素 $\tilde{c}_s \in \{0,1\}$ 被确定，即可完全刻画集群的宏观量测配置 。这一代数结构表明，异构编队的可观测性深度解耦为“底层网络几何刚性$\mathcal{R}_{\mathbf{p}}$”与“顶层传感节点配置$\tilde{C}$”的叠加效应，为后续基于矩阵分解的传感器优化部署提供了直接的理论抓手。
\section{基于弱可观测量测条件的传感器配置策略}

在确立了无人机编队系统的异构量测模型及可观测性矩阵后，本节旨在解决一个核心的工程问题：在保证整个集群系统状态弱可观测的前提下，如何以最小的传感器成本，确定编队中绝对位置传感器的分配策略，即确定能够满足系统可观测性的位置量测矩阵 $\tilde{C}$。

\subsection{基于矩阵QR分解技术的编队量测配置}
为确定满足系统可观测性的位置量测矩阵 $\tilde{C}$ 的主对角线元素，本节引入基于带主列索引的矩阵QR分解技术 。首先，给出带主列索引的 QR分解引理：

\textbf{引理 3.1} 给定矩阵 $A\in\mathbb{R}^{m\times n}$ 且 $m\ge n$，$A$ 的带主列索引的 QR 分解可以写为：
\begin{equation}
	Q^{\top}AP=R'=
	\begin{bmatrix}
		R_{11} & R_{12}\\
		\mathbf{0}_{(m-r)\times r} & \mathbf{0}_{(m-r)\times (n-r)}
	\end{bmatrix}
\end{equation}
其中 $Q \in \mathbb{R}^{m\times m}$ 为正交矩阵，$P=P_{1} P_{2}\cdots P_{r}$ 为置换矩阵，$R_{11}\in\mathbb{R}^{r\times r}$ 为非奇异的上三角矩阵，$R_{12}\in\mathbb{R}^{r\times (n-r)}$。置换矩阵 $P$ 由一个列置换索引向量 $\mathbf{p}(\cdot)$ 编码，即矩阵$ R' $的第$j$列对应矩阵$Q^{\top}A $的第$\mathbf{p}(j)$列。

利用引理3.1，对无人机编队的刚性矩阵 $\mathcal{R}_{\mathbf{p}}$ 展开分析。如果 $M \ge Nn$，其中$M$表示传感器网络边中的数量，$N$表示集群无人机数量，$n$表示无人机的位置维度，对其进行矩阵分解，得到：
\begin{equation}
	Q^{\top} \mathcal{R}_{\mathbf{p}} P 
	= \mathcal{R}_{\mathbf{p}}' 
	= 
	\begin{bmatrix}
		\mathcal{R}_{11} & \mathcal{R}_{12}\\
		\mathbf{0}_{(M-r) \times r} & \mathbf{0}_{(M-r) \times (Nn-r)}
	\end{bmatrix}.
	\label{刚性矩阵的分解式}
\end{equation} 
由于 $\mathcal{R}_{11}$ 是非奇异的，因此 $\mathcal{R}_{\mathbf{p}}'$ 的前 $r$ 列全为主元列。考虑到矩阵 $\mathcal{R}_{\mathbf{p}}'$ 是由 
$Q^{\top} \mathcal{R}_{\mathbf{p}}$ 经过列变换得到的，具体的变换顺序由列置换矩阵 $P$ 表示。因此，向量 $\mathbf{p}$ 可用于表示 $\mathcal{R}_{\mathbf{p}}$ 中所有主元列的索引。具体而言，$ \mathcal{R}_{\mathbf{p}}' $的第$j$列对应矩阵$Q^{\top} \mathcal{R}_{\mathbf{p}} $的第$\mathbf{p}(j)$列。

如果 $M < Nn$，由于行数小于列数，无法直接应用引理 3.1。因此需要构造增广矩阵 
$\bar {\mathcal{R}}_{\mathbf{p}} = 
[\mathcal{R}_{\mathbf{p}} ; \mathbf{0}_{(Nn-M) \times Nn}]$。然后，对该增广矩阵进行分解：
\begin{equation}
	\bar Q^{\top} \bar {\mathcal{R}}_{\mathbf{p}} \bar {P} 
	= \bar {\mathcal{R}}_{\mathbf{p}}'
	\label{刚性矩阵增广矩阵的分解式}
\end{equation}
同样可以识别出代表 $ \bar {\mathcal{R}}_{\mathbf{p}} $ 主元列的索引向量 $\bar{\mathbf{p}}$。根据 $\bar {\mathcal{R}}_{\mathbf{p}}$ 的构造结构，$\bar{\mathbf{p}}$同样对应于 ${\mathcal{R}}_{\mathbf{p}}$ 的主元列索引。

将 ${\mathcal{R}}_{\mathbf{p}}$ 的主元列索引集合记为：
\begin{equation}
	\mathcal{P}=\{i:\text{若 } M\ge Nn, i = \mathbf{p}(j), \text{否则 } i=\bar{\mathbf{p}}(j), j\in\mathcal{I}_{Nn}\},
\end{equation}
并将自由列索引集合记为：
\begin{equation}
	\bar{\mathcal{P}}=\complement_{\mathcal{I}_{Nn}}\mathcal{P}
\end{equation}

如前所述，确定矩阵$\tilde{C}  $的主对角元素即可完成位置传感器的配置。由可观测性矩阵的结构(\ref{可观测性矩阵})可知，${\mathcal{R}}_{\mathbf{p}}$ 的主元列也是 $ \mathcal{O}_{\mathbf{p}} $ 的主元列。对于一阶动力学(\ref{一阶动力学})，为了确保可观测性矩阵 $\mathcal{O}_{\mathbf{p}}$ 达到满列秩，即不存在自由列，可观测性矩阵 $\tilde{C}$ 的主对角线元素可按如下策略配置：
\begin{equation}
	\tilde{c}_{s} = 
	\begin{cases}
		1 & \text{如果 } s \in \bar{\mathcal{P}},\\
		0 & \text{否则 } ,
	\end{cases}
\end{equation}
该配置方程从数学上严密地规定了系统的状态输出要求。$\tilde{c}_{s}=1$ 的物理意义第 $\lceil s/n \rceil$ 个智能体的第 $(s \bmod n)$ 个状态维度是否需要被直接量测。这里的符号$\lceil / \rceil$表示除法向上取整，$ \bmod $表示整数除法取余数。值得注意的是，矩阵${\mathcal{R}}_{\mathbf{p}} $的分解式(\ref{刚性矩阵的分解式})与矩阵$ \bar {\mathcal{R}}_{\mathbf{p}}  $的分解式(\ref{刚性矩阵增广矩阵的分解式})并不是唯一的。因此$ \tilde{c}_i $的配置与量测也不是唯一的。

基于上述确定的主元列与传感器配置策略，可得到保证系统局部弱可观测性的充分条件。

\textbf{引理 3.2} 考虑一个无穷小刚性框架 $\mathcal{F}$的无人机编队系统(\ref{一阶动力学})。如果相对距离量测$ \bar {\mathbf{z}}_i $由传感器网络$\mathcal{G}$配置，位置量测$ \tilde{\mathbf{z}}_i$由位置量测矩阵$\tilde{C}_i $配置，则该无人机编队系统在 $\mathbf{p}^{*}$ 处是局部弱可观测的。

\textbf{证明：} 对无人机编队的刚性矩阵 ${\mathcal{R}}_{\mathbf{p}^*}$ 进行矩阵分解(\ref{刚性矩阵的分解式})可得矩阵的自由列索引集合$ \bar{\mathcal{P}} $ 与式：
\begin{equation}
	{\mathcal{R}}_{\mathbf{p}^*}=  Q \mathcal{R}_{\mathbf{p}^*}' P^{-1},
\end{equation}
结合可观测矩阵构造式(\ref{可观测性矩阵})可知：
\begin{equation}
	\mathcal{O}_{\mathbf{p}^*} = 
	\begin{bmatrix}
		\mathcal{R}_{\mathbf{p}^*}  \\
		\tilde{C}
	\end{bmatrix} = \begin{bmatrix}
	Q & \mathbf{0} \\
	\mathbf{0} & I_{Nn}
	\end{bmatrix}
	\begin{bmatrix}
		\mathcal{R}_{\mathbf{p}^*}'\\ 
		\tilde{C} P^{-1}
	\end{bmatrix}
	P^{-1}.
\end{equation}
根据列置换矩阵 $P$ 和对角矩阵 $\tilde{C}$ 的定义与配置，可以推导出：
\begin{equation}
	\tilde{C} P^{-1} = 
	\left[
		\mathbf{0}_{Nn \times \text{rank} (\mathcal{R}_{\mathbf{p}^*}')}
		\arrowvert
		\cdots \quad
		{\mathbf{1}_{Nn}}^1 \quad
		\cdots
	\right],
	\quad
	i \in \bar{\mathcal{P}},
\end{equation}
其中，$ \text{rank}(\mathcal{R}_{\mathbf{p}}') = \text{rank}(\mathcal{R}_{11}) $且$ \text{rank}([\cdots,{\mathbf{1}_{Nn}}^1,\cdots])=Nn-\text{rank}(\mathcal{R}_{\mathbf{p}^*}) $ 。
由分块对角矩阵的性质可得到$  \text{rank}(\mathcal{O}_{\mathbf{p}^*}) = \text{rank}(\mathcal{R}_{11})+ \text{rank}([\cdots,{\mathbf{1}_{Nn}}^1,\cdots]) = Nn $。可观测性矩阵列满秩，因此系统在 $\mathbf{p}^{*}$ 处满足局部弱可观测性。证明完毕。


\subsection{异构传感器配置实例分析}

为了更直观地展示上述基于QR分解的量测配置方法，本节考虑一个三维空间中由 $N=4$ 架无人机组成的编队实例。假设编队在全局坐标系下的期望位置分别为 $p_{1}^{*}=[1,1,1]^{\top}$，$p_{2}^{*}=[2,1,2]^{\top}$，$p_{3}^{*}=[1,2,2]^{\top}$，$p_{4}^{*}=[2,2,1]^{\top}$。无人机集群的通信拓扑为全连通图，即边集为 $\mathcal{E}=\{(i,j)\mid i,j\in\{1,2,3,4\}\}$，共包含 $M=6$ 条距离量测边。

在该几何构型下，系统的刚性矩阵秩为 $\text{rank}(\mathcal{R}_{p^{*}})=6$，满足无穷小刚性条件。由于 $M<Nn$ ($6 < 12$)，根据增广矩阵的QR分解方法，可以计算得到该刚性矩阵的主元列索引集合为 $\mathcal{P}=\{1,2,4,5,6,9\}$，从而确定其自由列索引集合为 $\overline{\mathcal{P}}=\{3,7,8,10,11,12\}$。

根据式(3-15)的传感器配置规则，为确保系统状态满秩可观测，必须对自由列对应的状态维度进行直接量测。由此可确定各无人机所需的位置量测矩阵 $\tilde{C}_{i}$：
\begin{equation}
	\tilde{C}_{1}=\text{diag}\{0,0,1\},\quad \tilde{C}_{2}=\text{diag}\{0,0,0\}
\end{equation}
\begin{equation}
	\tilde{C}_{3}=\text{diag}\{1,1,0\},\quad \tilde{C}_{4}=\text{diag}\{1,1,1\}
\end{equation}

在实际工程场景中，上述异构量测矩阵直接对应于不同无人机的传感器载荷配置方案：
\begin{itemize}
	\item \textbf{锚点无人机（无人机 4）：} 具备全维度位置量测（$\tilde{C}_{4}=I_{3}$），可搭载激光惯性里程计(LIO)等高级传感器，提供完整的三维全局坐标 $(p_{4x},p_{4y},p_{4z})$。
	\item \textbf{受限锚点无人机（无人机 1 和 3）：} 提供部分维度的绝对位置信息。例如，无人机 1 仅需搭载单线激光测距仪获取高度信息 $p_{1z}$；无人机 3 则可通过下视视觉相机进行特征点匹配，仅获取水平面坐标 $(p_{3x},p_{3y})$。
	\item \textbf{非锚点无人机（无人机 2）：} 无需搭载任何绝对位置传感器（$\tilde{C}_{2}=0_{3\times 3}$），仅依赖机间通信与 UWB 等测距模块获取相对距离量测即可实现协同定位。
\end{itemize}

通过上述合理的异构传感器配置，编队能够在大幅降低系统整体载荷与成本的前提下，维持全局坐标系下的弱可观测性，从而为后续的分布式状态估计提供充足的信息源。
